% Schematic domains and equatorial roots for a=-1, p in [p_0,P].
% The arrows in panel (a) denote positive-time transition maps.
% Panel (b) is an ordered root diagram, not an orbit on the equator.
\begin{figure}[htb]
\centering
\setlength{\unitlength}{1pt}
\begin{picture}(410,182)
 \put(205,172){\makebox(0,0){\small (a) The transition on fixed physical sections}}
 \put(45,148){\makebox(0,0){\(\Sigma_-\)}}
 \put(155,148){\makebox(0,0){\(X=-L\)}}
 \put(255,148){\makebox(0,0){\(X=L\)}}
 \put(365,148){\makebox(0,0){\(\Sigma_+\)}}
 \linethickness{.6pt}
 \put(142,128){\vector(-1,0){85}}
 \put(242,128){\vector(-1,0){74}}
 \put(353,128){\vector(-1,0){86}}
 \put(100,113){\makebox(0,0){\small outgoing tail}}
 \put(205,113){\makebox(0,0){\small compact passage}}
 \put(310,113){\makebox(0,0){\small incoming tail}}
 \put(45,97){\makebox(0,0){\(\dfrac{-\delta-\gamma}{s}\)}}
 \put(365,97){\makebox(0,0){\(\dfrac{\delta-\gamma}{s}\)}}
 \put(205,69){\makebox(0,0){\small (b) Equatorial roots in the same coordinate}}
 \linethickness{.35pt}
 \put(24,42){\vector(1,0){368}}
 \put(400,42){\makebox(0,0){\(X\)}}
 \put(195,39){\line(0,1){6}}
 \put(195,29){\makebox(0,0){\small \(0\)}}
 \put(74,42){\circle*{3}}
 \put(219,42){\circle*{3}}
 \put(338,42){\circle*{3}}
 \put(74,54){\makebox(0,0){\(-R\)}}
 \put(219,54){\makebox(0,0){\(x_+\)}}
 \put(338,54){\makebox(0,0){\(x_f\)}}
 \put(74,18){\makebox(0,0){\small \(-R\to-\infty\)}}
 \put(219,18){\makebox(0,0){\small \(x_+\to p^{-1}\)}}
 \put(338,18){\makebox(0,0){\small \(x_f\to+\infty\)}}
 \put(74,0){\makebox(0,0){\small \(\rho_-\to0\)}}
 \put(219,0){\makebox(0,0){\small \(\rho_0\le\rho_+\le\rho_1\)}}
\end{picture}
\caption{Two consequences of the rescaling \(v=\gamma+sX\),
\(z=s^2Y\), for \(\alpha_0=-s^2\) and \(p\in[p_0,P]\).
The positions and distances are schematic. In (a), a fixed compact
core does not include the tails leading to \(\Sigma_\pm\);
the arrows denote the three factors of the positive-time transition.
In (b), the ordered zeros of \(M\) are shown as \(s,\gamma\to0\).
The negative incident saddle escapes and its characteristic number
vanishes, although the positive incident saddle remains on a compact
scale. The axis arrow gives the orientation of \(X\), not the flow.}
\label{fig:tails-roots}
\end{figure}
